% Paper B: CMKL - Continual Multimodal KG Learner
% Target venues: ICML, NeurIPS, KDD, WWW

\documentclass{article}

% TODO: Replace with venue-specific style (e.g., \usepackage{neurips_2026})
\usepackage[utf8]{inputenc}
\usepackage{amsmath,amssymb}
\usepackage{graphicx}
\usepackage{booktabs}
\usepackage{hyperref}
\usepackage{natbib}
\usepackage{algorithm}
\usepackage{algorithmic}

\title{CMKL: Modality-Aware Continual Learning \\
for Evolving Biomedical Knowledge Graphs}

\author{
  Author Names \\
  KAUST \\
  \texttt{email@kaust.edu.sa}
}

\begin{document}

\maketitle

\begin{abstract}
% TODO: Write abstract after experiments are complete
Abstract placeholder. We propose CMKL, the first method to address
modality-specific forgetting in continual knowledge graph learning.
CMKL uses modality-aware regularization and multimodal memory replay
to leverage complementary information across structural, textual,
and molecular modalities.
\end{abstract}

% Section 1: Introduction (1 page)
\section{Introduction}
\label{sec:intro}
% Core challenge of continual multimodal biomedical KG reasoning
% Gap: no existing method handles modality-specific forgetting

% Section 2: Related Work (1 page)
\section{Related Work}
\label{sec:related}
% CKGE methods (EWC, LKGE, BER)
% Multimodal graph learning (MSCGL, MoDE)
% Biomedical KG reasoning

% Section 3: Problem Formulation (0.5 page)
\section{Problem Formulation}
\label{sec:problem}
% Formal definition of continual multimodal KG learning

% Section 4: Method (2.5 pages)
\section{CMKL: Continual Multimodal KG Learner}
\label{sec:method}

\subsection{Modality-Specific Encoders}
% R-GCN (structural), BiomedBERT (textual), Morgan FP MLP (molecular)

\subsection{Cross-Modal Attention Fusion}
% Bidirectional cross-attention, residual connections

\subsection{Modality-Aware EWC}
% Per-modality Fisher matrices, per-modality lambda values

\subsection{Multimodal Memory Replay}
% K-means diverse selection, embedding storage

% Section 5: Experiments (2 pages)
\section{Experiments}
\label{sec:experiments}
% Setup, main results vs. baselines, ablation studies

% Section 6: Analysis (1 page)
\section{Analysis}
\label{sec:analysis}
% Inter-modal vs. intra-modal forgetting
% Modality complementarity analysis
% Per-modality Fisher visualization

% Section 7: Conclusion (0.5 page)
\section{Conclusion}
\label{sec:conclusion}
% Summary, limitations, future work

\bibliographystyle{plainnat}
\bibliography{references}

\end{document}
